% Relativistic Chapter XII, §107–109: Gravitational Instability of Infinite Cylinder
% Relativistic generalization of Chandrasekhar, Chapter XII §107–109
% Conventions: see RELATIVISTIC_CONVENTIONS.md

\chapter{The Stability of Relativistic Jets and Cylinders}
\label{ch:rel-12}

%% ================================================================
\section{Introduction: relativistic jets and cylinders}\label{sec:rel-12-107}

The non-relativistic treatment of Chapter~XII opens with the classical
problem of the instability of a cylindrical jet, tracing its origins to
the work of Plateau and Lord Rayleigh. In the non-relativistic theory
the gravitational instability of an infinite self-gravitating cylinder is
governed by the Poisson equation and instantaneous Newtonian gravity;
the mode of maximum instability determines the characteristic
fragmentation scale.

In the relativistic regime these considerations must be revisited for
at least three reasons:
\begin{enumerate}
\item \textbf{Active galactic nuclei (AGN) jets.}  Relativistic jets
  launched from supermassive black holes propagate as highly collimated
  cylinders with bulk Lorentz factors $\Lf\sim 10$--$50$.  Their
  stability against varicose (pinch) and kink modes determines jet
  morphology and energy transport over kiloparsec to megaparsec scales.
  The internal energy density is comparable to the rest-mass energy
  density, so that the relativistic enthalpy $\enthalpy = (\varepsilon +
  p)/c^{2}$ replaces the Newtonian density $\rho_{0}$ in the inertia of
  the fluid.

\item \textbf{Cosmic strings.}  Topological defects formed during
  symmetry-breaking phase transitions in the early universe are
  effectively infinite relativistic cylinders of energy.  Their
  gravitational self-interaction is intrinsically general-relativistic:
  the deficit angle of the surrounding spacetime replaces the Newtonian
  gravitational potential, and perturbations propagate at the speed of
  light rather than instantaneously.

\item \textbf{Gamma-ray burst (GRB) jets and neutron-star jets.}
  The ultra-relativistic outflows from compact-object mergers and
  collapsars are pressure-supported cylinders in which $p/\varepsilon$
  can be of order unity.  Gravitational fragmentation of such jets
  competes with magnetic and hydrodynamic instabilities in shaping the
  observed variability.
\end{enumerate}

\paragraph{Key physical differences from the Newtonian theory.}
\begin{itemize}
\item \emph{Finite signal speed.}  In Newtonian gravity, gravitational
  perturbations propagate instantaneously; the Poisson equation
  $\nabla^{2}\mathfrak{B} = -4\pi G\rho$ is elliptic.  In general
  relativity, gravitational perturbations propagate at the speed of
  light; the linearized Einstein equations are hyperbolic.  This
  introduces a \emph{retardation} that fundamentally modifies the
  dispersion relation at wavelengths comparable to $c/\sqrt{G\rho}$.

\item \emph{Inertia of energy and pressure.}  The source of gravity is
  the energy-momentum tensor $\emt$, not the mass density alone.  For a
  perfect fluid, both $\varepsilon$ and $p$ contribute to self-gravity,
  so that pressure---which is stabilizing in the Newtonian case---has a
  partially \emph{destabilizing} role in general relativity (cf.\ the
  Tolman--Oppenheimer--Volkoff effect).

\item \emph{Causal dissipation.}  When viscosity is present, the
  Navier--Stokes equation must be replaced by the Israel--Stewart theory
  to preserve causality; the resulting relaxation terms modify the
  viscous damping at short wavelengths.
\end{itemize}

Throughout this chapter we work in cylindrical coordinates
$(t,\varpi,\varphi,z)$ with the metric signature $(-,+,+,+)$.  The
speed of light $c$ and the gravitational constant $G$ are kept explicit.
The four-velocity satisfies $u^{\mu}u_{\mu}=-c^{2}$.  In the
post-Newtonian treatment of \S\ref{sec:rel-12-108}, we expand in the
dimensionless compactness parameter
\begin{equation}
\label{eq:rel-12-1}
\mathcal{C} \;\equiv\; \frac{G M_{\mathrm{lin}}}{c^{2} R}
  \;=\; \frac{\pi G \rho_{0} R^{2}}{c^{2}},
\end{equation}
where $M_{\mathrm{lin}} = \pi R^{2}\rho_{0}$ is the mass per unit
length and $R$ the cylinder radius.  The Newtonian limit is recovered as
$\mathcal{C}\to 0$.

%% ================================================================
\section{Gravitational instability of a relativistic infinite cylinder}
\label{sec:rel-12-108}

\subsection{Equilibrium configuration}
\label{sec:rel-12-108a}

Consider an infinite cylinder of radius $R$ composed of a perfect fluid
with constant rest-mass density $\rho_{0}$, energy density
$\varepsilon$, and pressure $p$.  In the non-relativistic limit,
$\varepsilon \to \rho_{0}c^{2}$ and the equilibrium pressure is
$p_{0}(\varpi) = \pi G \rho_{0}^{2}(R^{2}-\varpi^{2})$ [cf.\
equation~(22) of Chapter~XII].  In the relativistic case the
equilibrium is determined by the relativistic Euler equation
\begin{equation}
\label{eq:rel-12-2}
\frac{\dd p_{0}}{\dd \varpi} = -\frac{(\varepsilon_{0}+p_{0})}{c^{2}}
  \,\frac{\dd \Phi_{0}}{\dd \varpi},
\end{equation}
where $\Phi_{0}$ is the gravitational potential in the weak-field
(post-Newtonian) approximation.  For a uniform-density cylinder,
$\Phi_{0}(\varpi) = -\pi G \rho_{0}\varpi^{2}$ inside, and
equation~\eqref{eq:rel-12-2} yields
\begin{equation}
\label{eq:rel-12-3}
p_{0}(\varpi) = \pi G \rho_{0}^{2}(R^{2}-\varpi^{2})
  \left[1 + \frac{(\varepsilon_{0}+p_{0})}{2\rho_{0}c^{2}}\right]^{-1}
  \!\!.
\end{equation}
To first post-Newtonian (1PN) order, writing
$\varepsilon_{0}=\rho_{0}c^{2}(1+\Pi_{\mathrm{int}}/c^{2})$ where
$\Pi_{\mathrm{int}}$ is the specific internal energy, we obtain
\begin{equation}
\label{eq:rel-12-4}
p_{0}(\varpi) \;=\; \pi G\rho_{0}^{2}(R^{2}-\varpi^{2})
  \left[1 - \frac{\Pi_{\mathrm{int}}}{c^{2}}
  - \frac{p_{0}}{2\rho_{0}c^{2}} + O(\mathcal{C}^{2})\right].
\end{equation}
In the Newtonian limit ($c\to\infty$), the correction terms vanish and
the classical result is recovered.

\subsection{Linearized perturbation equations: post-Newtonian gravity}
\label{sec:rel-12-108b}

We consider axisymmetric perturbations of the form
\begin{equation}
\label{eq:rel-12-5}
\varpi = R + \epsilon_{0}\, e^{i(kz + m\varphi) + \sigma t},
\end{equation}
with $m=0$ for the gravitationally unstable varicose modes.

\paragraph{Replacement of the Poisson equation.}
In the Newtonian theory the gravitational potential perturbation
$\delta\mathfrak{B}$ satisfies the Poisson equation
$\nabla^{2}\delta\mathfrak{B} = -4\pi G\delta\rho$.  In the
post-Newtonian framework the relevant equation is obtained from the
linearized $(0,0)$-component of the Einstein field equations in the
weak-field limit:
\begin{equation}
\label{eq:rel-12-6}
\nabla^{2}\delta\Phi = -4\pi G\left(\delta\rho
  + \frac{\delta\varepsilon + 3\,\delta p}{2c^{2}}
  + \frac{2(\varepsilon_{0}+p_{0})\,\delta v_{i} v^{i}}{c^{4}}\right)
  + \frac{1}{c^{2}}\frac{\partial^{2}\delta\Phi}{\partial t^{2}}.
\end{equation}
The last term on the right, the wave operator, encodes the
\emph{finite propagation speed} of gravitational perturbations.  In the
Newtonian limit ($c\to\infty$) it vanishes and the standard Poisson
equation is recovered.

For modes with time-dependence $e^{\sigma t}$, the wave term contributes
$\sigma^{2}\delta\Phi/c^{2}$.  For the gravitationally unstable modes
($\sigma^{2}>0$, $\sigma$ real), we write
\begin{equation}
\label{eq:rel-12-7}
\left(\nabla^{2} - \frac{\sigma^{2}}{c^{2}}\right)\delta\Phi
  = -4\pi G\left(\delta\rho + \frac{\delta\varepsilon + 3\,\delta p}{2c^{2}}\right).
\end{equation}
This is a Helmholtz equation rather than the Laplace equation; the
solutions involve modified Bessel functions with argument
$\tilde{k}\varpi$ where
\begin{equation}
\label{eq:rel-12-8}
\tilde{k}^{2} = k^{2} + \frac{\sigma^{2}}{c^{2}}.
\end{equation}

\paragraph{Modified Bessel-function solutions.}
For an incompressible fluid ($\delta\rho = 0$ inside the cylinder), the
interior and exterior potentials become
\begin{equation}
\label{eq:rel-12-9}
\delta\Phi_{\mathrm{int}} = B\, I_{0}(\tilde{k}\varpi)\,e^{ikz+\sigma t},
\qquad
\delta\Phi_{\mathrm{ext}} = A\, K_{0}(\tilde{k}\varpi)\,e^{ikz+\sigma t},
\end{equation}
where $I_{0}$ and $K_{0}$ are modified Bessel functions.  The matching
conditions at $\varpi=R$ (continuity of $\delta\Phi$ and of its normal
derivative minus the surface term) yield, in direct analogy with the
Newtonian case [cf.\ equations~(9)--(12) of Chapter~XII],
\begin{equation}
\label{eq:rel-12-10}
A = 4\pi G\rho_{0}^{*}\,R\epsilon_{0}\,I_{0}(\tilde{x}),
\qquad
B = 4\pi G\rho_{0}^{*}\,R\epsilon_{0}\,K_{0}(\tilde{x}),
\end{equation}
where $\tilde{x}=\tilde{k}R$ and
\begin{equation}
\label{eq:rel-12-11}
\rho_{0}^{*} \;=\; \rho_{0}\!\left(1 + \frac{\varepsilon_{0}+3p_{0}}{2\rho_{0}c^{2}}\right)
\end{equation}
is the \emph{effective gravitating density} that accounts for the
contribution of energy and pressure to the source of gravity.

\subsection{Relativistic dispersion relation}
\label{sec:rel-12-108c}

Proceeding as in the Newtonian analysis---matching the radial velocity
to the deformed boundary and imposing the vanishing of the total stress
on the free surface---we arrive at the \textbf{relativistic dispersion
relation for varicose modes} ($m=0$):
\begin{equation}
\label{eq:rel-12-12}
\boxed{
\frac{\sigma^{2}}{4\pi G\rho_{0}^{*}}
  = \frac{\tilde{x}\,I_{1}(\tilde{x})}{I_{0}(\tilde{x})}
  \left[K_{0}(\tilde{x})\,I_{0}(\tilde{x}) - \frac{1}{2}\right],}
\end{equation}
where $\tilde{x} = R\sqrt{k^{2}+\sigma^{2}/c^{2}}$ and $\rho_{0}^{*}$
is defined in \eqref{eq:rel-12-11}.  This is an \emph{implicit}
equation for $\sigma$ (since $\tilde{x}$ depends on $\sigma$); it must
be solved self-consistently.

\paragraph{Newtonian limit.}
When $c\to\infty$, we have $\tilde{x}\to x = kR$, $\rho_{0}^{*}\to\rho_{0}$,
and equation~\eqref{eq:rel-12-12} reduces to the classical
Chandrasekhar dispersion relation [equation~(29) of Chapter~XII]:
\begin{equation}
\label{eq:rel-12-13}
\frac{\sigma_{0}^{2}}{4\pi G\rho_{0}}
  = \frac{x\,I_{1}(x)}{I_{0}(x)}
  \left[K_{0}(x)\,I_{0}(x) - \tfrac{1}{2}\right].
\end{equation}

\paragraph{Relativistic corrections to the characteristic frequencies.}
Expanding equation~\eqref{eq:rel-12-12} to first post-Newtonian order,
we write $\sigma^{2} = \sigma_{0}^{2}(1+\delta_{\mathrm{rel}})$ with
$\sigma_{0}^{2}$ the Newtonian growth rate.  A straightforward expansion
yields
\begin{equation}
\label{eq:rel-12-14}
\delta_{\mathrm{rel}} = \frac{\varepsilon_{0}+3p_{0}}{2\rho_{0}c^{2}}
  - \frac{\sigma_{0}^{2}R^{2}}{c^{2}}\,
  \frac{\partial}{\partial x^{2}}
  \ln\!\left[\frac{x\,I_{1}(x)}{I_{0}(x)}
  \Bigl(K_{0}(x)I_{0}(x)-\tfrac{1}{2}\Bigr)\right]
  + O(\mathcal{C}^{2}).
\end{equation}
The first term reflects the enhanced gravitating mass (destabilizing);
the second is the retardation correction from the finite speed of
gravity (stabilizing for the growing modes).  Both corrections are of
order $\mathcal{C}=\pi G\rho_{0}R^{2}/c^{2}$.

\paragraph{Critical wavenumber.}
The marginal wavenumber $\tilde{x}_{a}$ at which $\sigma^{2}=0$ is
determined by
\begin{equation}
\label{eq:rel-12-15}
I_{0}(\tilde{x}_{a})\,K_{0}(\tilde{x}_{a}) = \tfrac{1}{2}.
\end{equation}
Since $\sigma=0$ implies $\tilde{x}_{a}=x_{a}=kR$, the critical
wavenumber is \emph{unchanged} by the relativistic corrections:
\begin{equation}
\label{eq:rel-12-16}
x_{a} = 1.0668 \quad (\text{same as Newtonian}).
\end{equation}
This is physically transparent: at marginal stability, perturbations are
static and the wave-operator term in \eqref{eq:rel-12-7} vanishes
identically.  The critical \emph{wavelength}
$\lambda_{a}=2\pi R/x_{a}$ is therefore universal, independent of
compactness.

\paragraph{Mode of maximum instability.}
In contrast to the critical wavenumber, the wavenumber of maximum
instability \emph{is} shifted.  Let $x_{\max}^{\mathrm{(N)}}=0.580$ be
the Newtonian value.  To 1PN order,
\begin{equation}
\label{eq:rel-12-17}
x_{\max} = x_{\max}^{\mathrm{(N)}}
  - \alpha\,\frac{\pi G\rho_{0}R^{2}}{c^{2}} + O(\mathcal{C}^{2}),
\end{equation}
where $\alpha > 0$ is a numerical constant obtained by differentiating
the implicit dispersion relation.  The shift is towards \emph{longer
wavelengths}: relativistic retardation preferentially suppresses
short-wavelength modes.

\paragraph{Characteristic time for break-up.}
The maximum growth rate including 1PN corrections is
\begin{equation}
\label{eq:rel-12-18}
\sigma_{\max} \approx 0.2455\sqrt{4\pi G\rho_{0}}
  \left[1 + \frac{\varepsilon_{0}+3p_{0}}{4\rho_{0}c^{2}}
  - \beta\,\frac{\pi G\rho_{0}R^{2}}{c^{2}}\right],
\end{equation}
where $\beta>0$ encodes the retardation effect.  The enhanced
gravitating mass slightly \emph{increases} the growth rate while the
finite signal speed slightly \emph{decreases} it; which effect dominates
depends on the equation of state.

The characteristic time for fragmentation is
\begin{equation}
\label{eq:rel-12-19}
\tau_{\mathrm{rel}} = \frac{1}{\sigma_{\max}}
  \approx \frac{1}{0.2455\sqrt{4\pi G\rho_{0}}}
  \left[1 - \frac{\varepsilon_{0}+3p_{0}}{4\rho_{0}c^{2}}
  + \beta\,\frac{\pi G\rho_{0}R^{2}}{c^{2}}\right].
\end{equation}

\paragraph{Causality constraint.}
A fundamental requirement of the relativistic theory is that the growth
rate $\sigma$ cannot exceed $c/R$, i.e.\ no information can propagate
faster than light across the cylinder.  From
equation~\eqref{eq:rel-12-12}, one verifies that
\begin{equation}
\label{eq:rel-12-20}
\sigma^{2} \;\leq\; 4\pi G\rho_{0}^{*} \;<\; \frac{c^{2}}{R^{2}}
\quad \text{when } \mathcal{C}\ll 1,
\end{equation}
so the causality bound is automatically satisfied in the post-Newtonian
regime.  For ultracompact cylinders ($\mathcal{C}\sim 1$), the full
nonlinear Einstein equations must be used and the instability
saturates at a growth rate of order $c/R$.

\subsection{Stability to non-axisymmetric perturbations}
\label{sec:rel-12-108d}

For $m\neq 0$, the Newtonian result $\sigma_{m}^{2}<0$ (stability)
carries over to the relativistic case.  The proof is identical: the
inequality $K_{m}(\tilde{x})I_{m}(\tilde{x})<1/2$ holds for all
$m\neq 0$ and all $\tilde{x}>0$, so
\begin{equation}
\label{eq:rel-12-21}
\sigma_{m}^{2} < 0 \quad \text{for all } m\neq 0.
\end{equation}
Thus, as in the Newtonian theory, the relativistic cylinder is
\emph{stable to all purely non-axisymmetric perturbations}.

\subsection{Energy principle in general relativity}
\label{sec:rel-12-108e}

The alternative derivation of the dispersion relation via the energy
principle (Chapter~XII, \S108(a)) can be extended to the relativistic
setting.  In general relativity, the relevant energy functional is the
\emph{canonical energy} of Friedman and Schutz, defined for a
Lagrangian displacement $\boldsymbol{\xi}^{\mu}$ of the fluid:
\begin{equation}
\label{eq:rel-12-22}
\mathcal{E}_{c}[\boldsymbol{\xi}]
  = \delta^{2}\!\left(\int T^{0\mu}\,\dd\Sigma_{\mu}\right),
\end{equation}
where $\delta^{2}$ denotes the second-order perturbation and the
integral is over a spacelike hypersurface.

For a static equilibrium (no background flow), the canonical energy
reduces to the sum of kinetic and potential contributions:
\begin{equation}
\label{eq:rel-12-23}
\mathcal{E}_{c} = \underbrace{\delta^{2}\mathfrak{T}_{\mathrm{rel}}}_{\text{kinetic, }>0}
  \;+\; \underbrace{\delta^{2}\mathfrak{W}_{\mathrm{rel}}}_{\text{gravitational}}.
\end{equation}
The system is unstable if and only if
$\delta^{2}\mathfrak{W}_{\mathrm{rel}} < 0$ for some trial
displacement.

For an irrotational, incompressible displacement
$\boldsymbol{\xi} = \operatorname{grad}\phi$ with
$\phi = [\epsilon/(kI_{0}'(\tilde{x}))]\,I_{0}(\tilde{k}\varpi)\cos kz$,
one finds
\begin{equation}
\label{eq:rel-12-24}
\delta^{2}\mathfrak{W}_{\mathrm{rel}}
  = -2\pi^{2}G\rho_{0}^{*2}\,R^{2}\,
  \left[K_{0}(\tilde{x})\,I_{0}(\tilde{x}) - \tfrac{1}{2}\right]\epsilon^{2},
\end{equation}
which is negative for $\tilde{x}<x_{a}=1.0668$, confirming the
instability found from the dispersion relation.

The Lagrangian for the reduced degree of freedom $\epsilon(t)$ is
\begin{equation}
\label{eq:rel-12-25}
\mathcal{L}_{\mathrm{rel}} = \tfrac{1}{2}\pi R^{2}\enthalpy
  \left\{\frac{I_{0}(\tilde{x})}{\tilde{x}\,I_{1}(\tilde{x})}
  \,\dot{\epsilon}^{2}
  + 4\pi G\rho_{0}^{*}
  \left[K_{0}(\tilde{x})\,I_{0}(\tilde{x})-\tfrac{1}{2}\right]
  \epsilon^{2}\right\},
\end{equation}
where $\enthalpy = (\varepsilon_{0}+p_{0})/c^{2}$ is the relativistic
enthalpy density replacing the Newtonian $\rho_{0}$.  The equation of
motion $\ddot{\epsilon}=\sigma^{2}\epsilon$ reproduces
\eqref{eq:rel-12-12}.

%% ================================================================
\section{The effect of viscosity: Israel--Stewart theory on cylindrical
  instability}\label{sec:rel-12-109}

\subsection{Motivation: causal viscous transport}
\label{sec:rel-12-109a}

In Chapter~XII, \S109, the effect of Newtonian (Navier--Stokes) viscosity
on the gravitational instability was analysed.  The Newtonian viscous
stress is instantaneous: any shear is felt throughout the cylinder
without delay.  This is physically acceptable when the speed of viscous
signal propagation greatly exceeds all other relevant speeds, but
violates causality in the relativistic regime.

The Israel--Stewart theory of relativistic dissipative fluids remedies
this by introducing relaxation equations for the viscous stress tensor.
In the notation of the conventions document, the shear stress satisfies
\begin{equation}
\label{eq:rel-12-26}
\taupi\, u^{\alpha}\covd_{\!\alpha}\,\pi^{\langle\mu\nu\rangle}
  + \pi^{\mu\nu} = 2\shearvisc\,\sigma^{\mu\nu},
\end{equation}
where $\taupi$ is the shear relaxation time, $\shearvisc$ is the shear
viscosity coefficient, and $\sigma^{\mu\nu}$ is the shear tensor of the
flow.  The relaxation time $\taupi > 0$ ensures that viscous signals
propagate at the finite speed
\begin{equation}
\label{eq:rel-12-27}
v_{\mathrm{visc}} = \sqrt{\frac{\shearvisc}{\enthalpy\,\taupi}} \;<\; c.
\end{equation}
The Navier--Stokes limit is obtained as $\taupi\to 0$ with
$\shearvisc$ fixed.

\subsection{Linearized viscous perturbation equations}
\label{sec:rel-12-109b}

For axisymmetric perturbations ($m=0$) of the form $e^{ikz+\sigma t}$,
the relativistic momentum equation replacing~(59) of Chapter~XII becomes
\begin{equation}
\label{eq:rel-12-28}
\enthalpy\,\sigma\,u_{i}
  = -\partial_{i}\,\delta\!\left(\Phi + \frac{p}{\enthalpy c^{2}}\right)
  + \partial_{j}\,\delta\pi^{ij},
\end{equation}
where $\delta\pi^{ij}$ is the perturbed shear stress.  In Fourier space,
the Israel--Stewart equation~\eqref{eq:rel-12-26} becomes
\begin{equation}
\label{eq:rel-12-29}
(1 + \taupi\sigma)\,\delta\pi^{ij}
  = 2\shearvisc\,\delta\sigma^{ij},
\end{equation}
so that the \emph{effective viscosity} for perturbations of growth rate
$\sigma$ is
\begin{equation}
\label{eq:rel-12-30}
\boxed{\nu_{\mathrm{eff}}(\sigma)
  = \frac{\nu}{1+\taupi\sigma},}
\end{equation}
where $\nu = \shearvisc/\enthalpy$ is the kinematic viscosity.  This is
the central simplification: the Israel--Stewart relaxation reduces to a
frequency-dependent effective viscosity that interpolates between the
Navier--Stokes value ($\taupi\sigma\ll 1$) and vanishing viscosity
($\taupi\sigma\gg 1$).

\subsection{The case when viscous effects are dominant}
\label{sec:rel-12-109c}

In the limit of large viscosity ($\nu\to\infty$ with $\taupi$ fixed),
the Newtonian analysis of \S109(a) gave
\begin{equation}
\label{eq:rel-12-31}
\sigma^{(\mathrm{N})} = \frac{4\pi G\rho_{0}R^{2}}{\nu}\,
  \frac{K_{0}(x)I_{0}(x)-\frac{1}{2}}{x^{2}
  [I_{1}^{2}(x)/I_{0}^{2}(x)-(1+x^{2})^{-1}]}.
\end{equation}
In the relativistic generalization, two modifications enter:
\begin{enumerate}
\item The Newtonian density $\rho_{0}$ is replaced by the effective
  gravitating density $\rho_{0}^{*}$ defined in \eqref{eq:rel-12-11}.
\item The kinematic viscosity $\nu$ is replaced by the effective
  viscosity $\nu_{\mathrm{eff}}(\sigma)$ of \eqref{eq:rel-12-30}.
\end{enumerate}
The resulting implicit equation for $\sigma$ in the high-viscosity limit is
\begin{equation}
\label{eq:rel-12-32}
\sigma\,(1+\taupi\sigma) = \frac{4\pi G\rho_{0}^{*}\,R^{2}}{\nu}\,
  \frac{K_{0}(x)I_{0}(x)-\frac{1}{2}}{x^{2}
  [I_{1}^{2}(x)/I_{0}^{2}(x)-(1+x^{2})^{-1}]}.
\end{equation}
Denoting the right-hand side by $S(x)$, this is a quadratic in $\sigma$:
\begin{equation}
\label{eq:rel-12-33}
\taupi\,\sigma^{2} + \sigma - S(x) = 0,
\end{equation}
with the physical root
\begin{equation}
\label{eq:rel-12-34}
\sigma = \frac{-1+\sqrt{1+4\taupi S(x)}}{2\taupi}.
\end{equation}

\paragraph{Newtonian limit.}
As $\taupi\to 0$, $\sigma\to S(x)$, recovering equation~(96) of
Chapter~XII (with $\rho_{0}^{*}\to\rho_{0}$).

\paragraph{Large relaxation time.}
As $\taupi\to\infty$ (or for rapidly growing modes with
$\taupi\sigma\gg 1$), $\sigma\to\sqrt{S(x)/\taupi}$.  The growth rate
is reduced compared to the Newtonian case: the Israel--Stewart
relaxation acts as an effective reduction of viscous coupling.  In this
regime the growth rate transitions from the viscosity-dominated
$\sigma\propto 1/\nu$ scaling to the intermediate
$\sigma\propto 1/\sqrt{\nu\taupi}$ scaling.

\paragraph{Causality bound.}
The growth rate \eqref{eq:rel-12-34} is bounded:
\begin{equation}
\label{eq:rel-12-35}
\sigma < \frac{1}{2\taupi}\left(-1 + \sqrt{1+\frac{16\pi G\rho_{0}^{*}\taupi R^{2}}{\nu}}\right)
  \;<\; \frac{1}{\taupi}\sqrt{\frac{4\pi G\rho_{0}^{*}R^{2}}{\nu}}
  \;<\; \frac{c}{R},
\end{equation}
the last inequality following from $\nu/\taupi < c^{2}$
[cf.\ \eqref{eq:rel-12-27}] and $\mathcal{C}\ll 1$.  Thus causality is
preserved.

\paragraph{Monotone behaviour and fragmentation.}
As in the Newtonian high-viscosity limit, the growth rate
\eqref{eq:rel-12-34} is a monotone decreasing function of $x$ for
$0<x<x_{a}$; there is no finite mode of maximum instability.  The
Israel--Stewart relaxation does not change this qualitative feature:
the cylinder still tends to fragment at the longest available wavelength
rather than at the circumference.

\subsection{General case with relativistic corrections}
\label{sec:rel-12-109d}

For intermediate values of the viscosity, the full characteristic
equation is obtained by replacing $\nu\to\nu_{\mathrm{eff}}(\sigma)$
and $\rho_{0}\to\rho_{0}^{*}$ in the Newtonian equation~(86) of
Chapter~XII.  With $\ell^{2}=k^{2}+\sigma/\nu_{\mathrm{eff}}$, the
Israel--Stewart modification gives
\begin{equation}
\label{eq:rel-12-36}
\ell^{2} = k^{2} + \frac{\sigma(1+\taupi\sigma)}{\nu},
\end{equation}
so that $y = \ell R$ satisfies
\begin{equation}
\label{eq:rel-12-37}
\frac{\sigma R^{2}}{\nu}\,(1+\taupi\sigma) = y^{2}-x^{2}.
\end{equation}
Equation~(90) of Chapter~XII then becomes
\begin{equation}
\label{eq:rel-12-38}
(x^{2}+y^{2})\frac{I_{1}(x)}{I_{0}(x)}
  \left\{\frac{I_{0}(x)}{I_{0}(y)}\frac{I_{1}(y)}{I_{1}(x)}
  - \frac{2xy}{x^{2}+\tfrac{1}{2}y^{2}}\right\}
  [x^{2}-y^{2}]
  = -\frac{J^{*}}{x^{2}+y^{2}}\,
  \left[K_{0}(x)I_{0}(x)-\tfrac{1}{2}\right],
\end{equation}
where the \textbf{relativistic viscous parameter} is
\begin{equation}
\label{eq:rel-12-39}
J^{*} = \frac{4\pi G\rho_{0}^{*}R^{4}}{\nu^{2}}
  = J\,\frac{\rho_{0}^{*}}{\rho_{0}}
  = J\!\left(1+\frac{\varepsilon_{0}+3p_{0}}{2\rho_{0}c^{2}}\right),
\end{equation}
with $J = 4\pi G\rho_{0}R^{4}/\nu^{2}$ the Newtonian parameter
[cf.\ equation~(89) of Chapter~XII].

The system of equations~\eqref{eq:rel-12-37} and \eqref{eq:rel-12-38}
implicitly determines the growth rate $\sigma$ as a function of $x$,
$J^{*}$, and the Israel--Stewart parameter
\begin{equation}
\label{eq:rel-12-40}
\mathcal{T} \;\equiv\; \frac{\taupi\nu}{R^{2}}.
\end{equation}

\paragraph{Key results.}
\begin{enumerate}
\item \emph{Instability criterion unaffected.}  As in the Newtonian
  case, the modes $0<x<x_{a}$ are unstable and the modes $x>x_{a}$ are
  stable; viscosity (whether Navier--Stokes or Israel--Stewart) does not
  change the range of unstable wavenumbers.

\item \emph{Mode of maximum instability.}  For finite $J^{*}$ and
  $\mathcal{T}=0$ (Navier--Stokes limit), the mode of maximum
  instability shifts to smaller $x$ (longer wavelengths) as $J^{*}$
  decreases (increasing viscosity), exactly as in the Newtonian theory.
  For $\mathcal{T}>0$ (Israel--Stewart), the effective viscosity is
  reduced at high $\sigma$, causing the mode of maximum instability to
  shift \emph{back towards} the inviscid value $x=0.580$---the
  relaxation partially undoes the viscous suppression of
  short-wavelength modes.

\item \emph{Dispersion curves.}  For a given $(J^{*},\mathcal{T})$,
  the dispersion relation can be obtained numerically by solving
  equations~\eqref{eq:rel-12-37}--\eqref{eq:rel-12-38} for a sequence
  of $y$ values at each $x$.  The progressive shifting of
  $x_{\max}$ with decreasing $J^{*}$ (Fig.~126 of Chapter~XII) is
  modulated by $\mathcal{T}$: larger relaxation times compress the
  family of curves towards the inviscid limit.
\end{enumerate}

\paragraph{Summary of gravitational signal speed.}
In all cases treated in this section, the crucial physical difference
from the Newtonian theory is that \emph{gravitational signals propagate
at the speed of light, not instantaneously}.  This manifests in three
ways:
\begin{enumerate}
\item The Poisson equation is replaced by a wave equation
  [equation~\eqref{eq:rel-12-7}], introducing the modified wavenumber
  $\tilde{k}$ [equation~\eqref{eq:rel-12-8}].
\item The viscous stress propagates at the finite speed
  $v_{\mathrm{visc}} < c$
  [equation~\eqref{eq:rel-12-27}], preventing the viscous
  characteristic equation from admitting superluminal modes.
\item The growth rate is bounded: $\sigma < c/R$
  [equations~\eqref{eq:rel-12-20} and \eqref{eq:rel-12-35}], in
  contrast to the Newtonian theory where $\sigma$ can in principle be
  arbitrarily large for sufficiently dense cylinders.
\end{enumerate}
